% ===========================================================================
\chapter{Metodologia e Metadados da Prospecção Tecnológica}
\label{cap_metodologia_metadados}
% ===========================================================================

A metodologia adotada para esta busca de anterioridade seguiu os preceitos científicos de prospecção patentária e revisão sistemática do estado da técnica, executada no dia \textbf{02 de setembro de 2026}. 

Foram consultadas as principais bases de dados nacionais e internacionais de documentos patentários e literatura científica revisada por pares, cobrindo de forma exaustiva o espectro técnico da computação aplicada e o domínio pedagógico da segurança do paciente hospitalar.

% ===========================================================================
\section{Eixos Temáticos e Formulação Booleana de Busca}
\label{sec_eixos_tematicos}
% ===========================================================================

A formulação das expressões booleanas de busca estruturou-se na interseção analítica de três eixos temáticos fundamentais:
\begin{enumerate}
    \item \textbf{Eixo 1 -- Segurança do Paciente e Gestão de Riscos}: Termos controlados em inglês e português cobrindo eventos adversos e mitigação de falhas clínicas (\textit{``patient safety''}, \textit{``adverse events''}, \textit{``risk management''}, \textit{``segurança do paciente''}, \textit{``eventos adversos''}, \textit{``gerenciamento de riscos''});
    \item \textbf{Eixo 2 -- Engenharia de Software e Tecnologias de Informação}: Termos relacionados a soluções de software e painéis interativos (\textit{``system''}, \textit{``software''}, \textit{``dashboard''}, \textit{``application''}, \textit{``sistema''});
    \item \textbf{Eixo 3 -- Educação e Treinamento em Saúde}: Descritores direcionados a ferramentas pedagógicas e formação profissional (\textit{``educational tool''}, \textit{``training''}, \textit{``learning''}, \textit{``education''}, \textit{``teaching''}, \textit{``educação''}, \textit{``treinamento''}, \textit{``ensino''}).
\end{enumerate}

A conjunção lógica dos três eixos expressa-se pela relação booleana:
\begin{equation}
\label{eq_busca_booleana}
\text{Expressão Global} = (\text{Eixo 1}) \;\land\; (\text{Eixo 2}) \;\land\; (\text{Eixo 3})
\end{equation}

% ===========================================================================
\section{Levantamento Quantitativo e Tabela de Metadados}
\label{sec_tabela_metadados}
% ===========================================================================

A Tabela \ref{tab_metadados_busca} consolida as bases bibliográficas e patentárias consultadas, as tipologias de documentos analisados, as estratégias booleanas exatas aplicadas e os quantitativos brutos e filtrados obtidos na data base da pesquisa.

\begin{table}[!ht]
\centering
\caption[Metadados de Busca e Levantamento Quantitativo por Base de Dados]{Metadados de Busca e Levantamento Quantitativo por Base de Dados (02/09/2026)}
\label{tab_metadados_busca}
\footnotesize
\begin{tabularx}{\textwidth}{p{2.6cm} p{1.8cm} X p{2.2cm}}
\toprule
\textbf{Base de Dados} & \textbf{Tipologia} & \textbf{Expressão Booleana de Busca (String)} & \textbf{Resultados} \\ 
\midrule
\textbf{Scopus} & Artigos & \texttt{TITLE-ABS-KEY(("adverse events" OR "patient safety" OR "risk management") AND ("educational tool" OR "training" OR "learning" OR "education" OR "teaching") AND ("system" OR "software" OR "dashboard" OR "application"))} & 26.185 brutos \\ 
\addlinespace
\textbf{Web of Science} & Artigos & \texttt{Tópico: (("adverse events" OR "patient safety" OR "risk management") AND ("educational tool" OR "training" OR "learning" OR "education" OR "teaching") AND ("system" OR "software" OR "dashboard" OR "application"))} & 12.338 brutos \\ 
\addlinespace
\textbf{SciELO} & Artigos & \texttt{Tópico: (("adverse events" OR "patient safety" OR "risk management") AND ("educational tool" OR "training" OR "learning" OR "education" OR "teaching") AND ("system" OR "software" OR "dashboard" OR "application"))} & 283 brutos \\ 
\addlinespace
\textbf{WIPO Patentscope} & Patentes & \texttt{FP: ((("adverse events" OR "patient safety" OR "risk management") AND ("educational tool" OR "training" OR "learning" OR "education" OR "teaching") AND ("system" OR "software" OR "dashboard" OR "application")))} & 542 globais \newline (86 sob G16H) \\ 
\addlinespace
\textbf{Portal CAPES} & Artigos & \texttt{Anyfield: (("eventos adversos" OR "segurança do paciente" OR "gerenciamento de riscos") AND ("educação" OR "treinamento" OR "ensino") AND ("sistema" OR "software" OR "dashboard"))} & 292 brutos \newline (65 filtrados) \\ 
\bottomrule
\end{tabularx}
\fonte{Elaborado pelo autor (2026).}
\end{table}

% ===========================================================================
\section{Critérios de Triagem e Seleção do Estado da Técnica}
\label{sec_triagem_selecao}
% ===========================================================================

A triagem inicial demonstrou expressivo volume bruto nas bases indexadoras multidisciplinares internacionais (Scopus e Web of Science), demonstrando a relevância acadêmica do tema da segurança do paciente. No entanto, a quase totalidade desses estudos aborda simulações analógicas presenciais com manequins de ressuscitação ou análises teóricas sem artefatos de software desenvolvidos.

Para a delimitação do estado da técnica com real valor preditivo para o SIGEA-GTT, foram adotados os seguintes critérios de refinamento:
\begin{itemize}
    \item \textbf{Nas bases de patentes (WIPO Patentscope)}: Refinamento dos 542 documentos pelo código IPC \textbf{G16H} (\textit{Tecnologia da Informação e Comunicação aplicada à Saúde}), resultando em 86 documentos patentários, dos quais foram selecionados os quatro pedidos de invenção com maior proximidade arquitetural;
    \item \textbf{Nas bases científicas nacionais e regionais (Portal CAPES e SciELO)}: Aplicação de filtros por língua (português e inglês), atualidade (2020 a 2026) e presença de arquiteturas de software operacionais com validação empírica, resultando no subconjunto de 65 artigos relevantes, dentre os quais três destacam-se como paradigmas metodológicos centrais.
\end{itemize}
